\chapter{Introduction}
\label{ch:intro}

\section{Background and Motivation}

Dementia is a clinical syndrome in which memory and other cognitive abilities decline to the point that a person can no longer manage daily life independently. Alzheimer's disease is its most common cause. The underlying brain changes begin years before the diagnosis, and many people pass through an intermediate stage, mild cognitive impairment (MCI), in which deficits are measurable but daily function is preserved. Not everyone with MCI progresses, and some people who appear cognitively normal decline within a few years. Estimating who is likely to progress in the near future is useful for monitoring, for planning care and for selecting participants for clinical trials.

A wide range of information is available for this purpose. Demographic characteristics, in particular age, are strong risk factors. Cardiovascular and metabolic conditions such as hypertension and diabetes have been associated with later dementia in epidemiological studies \cite{livingston2020,sierra2020,hainsworth2024,cao2024}. The APOE $\varepsilon$4 allele is the major genetic risk factor for late-onset Alzheimer's disease \cite{wolfe2018}. Neuropsychological tests and clinical ratings measure the current state of cognition and function, and amyloid and tau PET imaging measures the pathology directly \cite{chandra2019,leuzy2019}. What is less clear is how much each source adds once the others are known, and whether the added complexity of modern machine-learning models is justified for such data.

\section{How the Research Question Developed}

The project began with a broad aim: to predict cognitive decline from cardiovascular risk factors combined with longitudinal cognitive data, APOE and PET biomarkers, and to explain the predictions with supporting literature. Before any model was built, the data were audited in detail. The audit showed that the clinical diagnosis and the Clinical Dementia Rating are recorded consistently at every visit across twenty years of data, that roughly one in three participants with MCI progresses to dementia within three years, and that most PET scans are too recent to observe long-term outcomes. On this basis the aim was narrowed to one precise and measurable question: the probability of a dementia diagnosis within three years of a clinic visit. The original cardiovascular hypothesis was retained, but as a question to be tested rather than an assumption: whether cardiovascular information adds predictive value once current cognition is known.

\section{Changes from the Initial Proposal}

The framework proposed at the zeroth review was designed before the data had been examined. Several of its components were changed once the data audit and the first experiments showed what the data could support. Table~\ref{tab:plan} records each change and its reason.

\begin{table}[H]
\centering\footnotesize
\caption{Planned components at the zeroth review and what was implemented.}
\label{tab:plan}
\begin{tabularx}{\textwidth}{@{}>{\raggedright\arraybackslash}p{3.3cm}>{\raggedright\arraybackslash}p{3.6cm}Y@{}}
\toprule
Planned & Implemented & Reason \\
\midrule
Target: cognitive decline and progression, not yet fixed & Dementia diagnosis within three years; CDR-SB rise of one point built as a secondary label & Diagnosis is recorded consistently at every visit in every form version, and the three-year horizon balances events against loss to follow-up \\
Dataset: NACC via the ADSP-PHC harmonised release, NHANES & NACC Freeze 74 (UDS and SCAN PET tables, raw PET images) & The UDS release contains all required clinical variables with longitudinal visits and links to the PET data; NHANES does not provide longitudinal dementia diagnoses \\
XGBoost or CatBoost for tabular data & Logistic regression, random forest, XGBoost, MLP and FT-Transformer & A wider comparison answers whether model complexity helps; the models performed equally, so CatBoost was not added \\
LSTM or Transformer for cognitive trajectories & Trajectory features and a GRU on landmark visits & A GRU is the smaller recurrent model suited to short, irregular sequences of at most six visits; history added little \\
3D ResNet for PET with attention-based fusion & Small 3D CNN for amyloid status; late fusion of PET measures; modality dropout for missing data & Only 5 of 446 participants with images converted to dementia, so an image-based fusion model for progression could not be trained reliably; an early prototype of this kind gave unstable results (Section~\ref{sec:prelim}) \\
SHAP, Grad-CAM, attention weights, DiCE & SHAP and Grad-CAM & Counterfactual explanations suggest what to change in a patient, which the observational data cannot support \\
RAG grounding of predictions & Evidence assistant with hybrid retrieval and guardrails & Implemented as planned, with retrieval evaluated against a fixed question set \\
Target AUROC 0.70 to 0.78 & AUROC 0.94 overall, 0.84 within MCI & The overall figure is higher than expected mainly because current cognitive status is a strong predictor; results are reported per diagnostic group for this reason \\
\bottomrule
\end{tabularx}
\end{table}

\subsection{Preliminary Experiment}
\label{sec:prelim}
An early prototype followed the original plan closely: a 3D CNN on raw PET volumes fused with clinical features through cross-attention, predicting a rise of at least 0.5 in CDR-SB after one year. Only 198 participants had both a usable image and a one-year outcome, of whom 17 progressed. Under five-fold cross-validation the fused model reached a mean AUROC of 0.65 with fold values ranging from 0.40 to 0.86, and tabular models alone did no better than 0.78 with a standard deviation above 0.1. With so few events the estimates were too unstable to support any conclusion. This motivated the change to the full clinical cohort, a three-year dementia outcome, and a separate, validated treatment of the PET images.

\section{Problem Statement}

Let $\mathbf{x}_{t_0}$ denote the information available about a person without dementia at an index visit $t_0$: demographics, cardiovascular and medical history, neuropsychological test results, the clinician's rating, genetics and, where available, PET measures. The problem is to estimate
\begin{equation}
\hat{p} = P\left(y = 1 \mid \mathbf{x}_{t_0}\right),
\end{equation}
where $y = 1$ if dementia is diagnosed within three years of $t_0$ and $y = 0$ if the person is followed for three years without such a diagnosis; to determine how much each category of information contributes to $\hat{p}$; and to do so without using any information that would not have been available at $t_0$.

\section{Research Questions and Objectives}

The study addresses six research questions:
\begin{enumerate}[label=\textbf{RQ\arabic*.},leftmargin=*]
\item How accurately can a dementia diagnosis within three years be predicted for people without dementia, and how does this differ between cognitively normal people and people with MCI?
\item Does cardiovascular and medical information improve the prediction beyond demographics, genetics and current cognition?
\item Do amyloid and tau PET measures improve the prediction?
\item Do deep-learning models outperform classical models on this kind of data?
\item Does information from earlier visits improve the prediction beyond the current visit?
\item Can a single model provide calibrated estimates when whole categories of information are missing?
\end{enumerate}

The corresponding objectives were to construct a leakage-free cohort and outcome from the NACC database; to harmonise neuropsychological tests across form versions; to train and compare five classifiers under a person-level split with held-out centres; to measure the contribution of each category by ablation; to explain predictions with Shapley values; to test the added value of PET and of earlier visits; to validate the model in subgroups, for calibration and in sensitivity analyses; to build and validate a PET image pipeline; to integrate the model into a research prototype with an evidence assistant; and to verify the integrity of the entire pipeline with independent checks.

\section{Contributions}

\begin{itemize}
\item A reproducible pipeline that turns the NACC database into a leakage-free three-year prediction task, with test harmonisation across form versions.
\item A controlled comparison of five model families showing that deep-learning models offer no advantage over logistic regression or gradient boosting on these data.
\item A quantified answer to the cardiovascular question: in this cohort and setting, cardiovascular and medical history adds about 0.001 AUROC once current cognition is known.
\item Results reported separately for cognitively normal participants and participants with MCI, and a measurement of how strongly the prediction depends on current cognitive status.
\item A PET image pipeline validated against the PET core's own measurements.
\item A research prototype with input plausibility checks and an evidence assistant restricted to published sources.
\item A set of independent verification checks that found and corrected two sources of information leakage.
\end{itemize}

\begin{figure}[H]
\centering
\includegraphics[width=\textwidth]{pipeline.png}
\caption{Overview of the work presented in this thesis.}
\label{fig:pipeline}
\end{figure}

\section{Organisation of the Thesis}

Chapter~\ref{ch:lit} reviews related work. Chapter~\ref{ch:data} describes the data. Chapter~\ref{ch:method} presents the prediction task, cohort, features, models and evaluation. Chapter~\ref{ch:petimages} describes the PET image pipeline. Chapter~\ref{ch:results} reports the results. Chapter~\ref{ch:system} describes the prototype. Chapter~\ref{ch:verification} describes the verification of the pipeline. Chapter~\ref{ch:discussion} discusses the findings and their limitations, and Chapter~\ref{ch:conclusion} concludes. The appendices list the features and the reproduction steps.
